\documentclass[10pt,a4paper]{amsart}
\usepackage[margin=19mm]{geometry}
\usepackage{amsmath,amssymb,mathtools}
\usepackage[scr=rsfs,cal=euler]{mathalfa}
\usepackage{xfrac}
\usepackage[unicode,hidelinks,bookmarks=false]{hyperref}
\newcommand{\ie}{i.e.\ }
\newcommand{\cf}{cf.\ }
\newcommand{\resp}{respectively\ }
\newcommand{\Subsection}{\subsection}
\providecommand{\nobreakdash}{\nobreak\hbox{-}}
\setlength{\emergencystretch}{2em}
\multlinegap0pt
\usepackage{amssymb}
\usepackage[shortlabels]{enumitem}
\usepackage{dsfont}
\usepackage{caption}
\usepackage{tikz}
\usepackage{xcolor}
\usepackage{mdframed}
\usepackage{csquotes}
\usepackage{float}
\usepackage{bm}
\usepackage{bbm}
\usepackage{mathtools}
\newtheorem{lemma}{Lemma}
\newtheorem{proposition}{Proposition}
\newcommand\E{\mathbb{E}}
\newcommand\R{\mathbb{R}}
\newcommand\dd{\mathrm{d}}
\newcommand\nut{\tilde{\nu}}
\newcommand\one{\mathds{1}}
\newcommand\x{\underline{x}}
\title{{\bf Tightness of Stationary Nodal Measures}}
\author{Louis Gass, Giovanni Peccati}
\date{Source: arXiv:2512.17524v1 (CC BY 4.0)}
\begin{document}
\maketitle

\noindent\emph{Source and licence.} {\bf Tightness of Stationary Nodal Measures}. Version arXiv:2512.17524v1 (CC BY 4.0), distributed by arXiv under the Creative Commons Attribution 4.0 International licence, http://creativecommons.org/licenses/by/4.0/, as recorded in the arXiv licence metadata for this exact version and shown on the abstract page https://arxiv.org/abs/2512.17524v1. Full licence notice: \texttt{licenses/arxiv-2512.17524-CC-BY-4.0.txt}.\par\smallskip

\noindent\textit{Editorial scope.} Counted unit: the article's proposition prop2 (source0.tex lines 831--834). The article's complete original proof of it is delivered with it (source0.tex lines 835--841, one \texttt{proof} environment of 7 lines). No mathematics is written, repaired or re-derived: the statement and the proof are the article's own lines, in the article's own notation, and no retained line is substituted or re-flowed.  Retained uncounted as the source-local context the proof needs: lines 727--736; lines 737--746; lines 762--767; lines 816--819.  Nothing was omitted from any retained span, so the package carries no omission record.  No retained line contains a citation, so the wrapper carries no bibliography and no bibliography entry.  No retained line contains a figure environment, an image-inclusion command or drawing code, so no image is delivered and the package declares no asset dependency.  Editorial formatting only, in addition: an AMS \texttt{amsart} preamble that restates the article's own declarations and macros used by the retained lines, namely the environments \texttt{lemma}, \texttt{proposition} on separate counters, together with the standard packages those lines need.  The retained environments print in this document as Lemma 1, Proposition 1 in order of appearance, whereas the article prints the same items with the numbering of its own full document.  The complete native source of the article as distributed in the arXiv e-print for this exact version is bundled unchanged as \texttt{sources/191321/source0.tex}.
\begin{lemma}
\label{lem1}
Let $f$ be a random field on $\R^d$ satisfying assumptions  $(H1)$ (resp. $(\widetilde{H}1)$). If $k<d$, then
\[E[\nut(A)^2] = \int_{A^2}F_2(\x)\dd \x,\]
\[E[\nut(A)\nut(B)] = \int_{A\times B} F_2(\x)\dd \x,\]
and
\begin{align*}
E[\nut(A)^2\nut(B)^2] = \int_{A^2\times B^2}F_4(\x)\dd x + 2\E[\nut(A)\nut(B)]^2 + \E[\nut(A)^2]\E[\nut(B)^2].
\end{align*}
\end{lemma}

\begin{lemma}
\label{lem2} Let the assumptions of Lemma \ref{lem1} prevail. If $k=d$, then
\[E[\nut(A)^2] = \int_{A^2}F_2(\x)\dd \x + \int_{A}F(x)\dd x,\]
\[E[\nut(A)\nut(B)] = \int_{A\times B} F_2(\x)\dd \x,\]
and
\begin{align*}
E[\nut(A)^2\nut(B)^2] &= \int_{A^2\times B^2}F_4(\x)\dd x + \int_{A\times B^2}F_3(\x)\dd x + \int_{A^2\times B}F_3(\x)\dd x + \int_{A\times B}F_2(\x)\dd x\\
&+\vphantom{\int^{\int}}\quad 2\E[\nut(A)\nut(B)]^2 + \E[\nut(A)^2]\E[\nut(B)^2].
\end{align*}
\end{lemma}

\begin{lemma}
\label{lem3}
Let $f$ be a random field on $\R^d$ satisfying assumptions  $(H1)$ (resp. $(\widetilde{H}1)$) and $(H2)$. Then there is a constant $C$, depending only on the process $f$, such that
\[|F_2(x,y)|\leq C\|x-y\|^{-\beta}\one_{\|x-y\|\leq 1} + Cg^2(\|x-y\|)\one_{\|x-y\|\geq 1}.\]
In particular, the function $F_2$ is integrable.
\end{lemma}

\begin{lemma}
\label{lem4}
 Let the assumptions of Lemma \ref{lem3} prevail. Then there is a constant $C$ depending only on the process $f$ such that
\[\int_{[0,t]^p} |F_p(x_1,\ldots,x_p)|\dd x_1,\ldots\dd x_p\leq C\left(\prod_{j=1}^i t_j\right)\left(\prod_{j=i+1}^d t_j\right)^{p/2}.\]\end{lemma}

\begin{proposition}
\label{prop2} Let the assumptions of Lemma \ref{lem1} prevail. Then there is a constant $C$ depending only on the process $f$ such that
\[\E[\nut(A)^2\nut(B)^2]\leq C\left(\prod_{j=1}^i t_j\right)\left(\prod_{j=i+1}^d t_j\right)^2.\]
\end{proposition}

\begin{proof}
Using the expression of the left hand side given by Lemma \ref{lem1} and Lemma \ref{lem2}, and the estimate given by Lemma \ref{lem4}, one obtains the bound
\begin{align*}
\E[\nut(A)^2\nut(B)^2]&\leq C\left(\prod_{j=1}^i t_j\right)\left[\sum_{p=2}^4 \left(\prod_{j=i+1}^d t_j\right)^{p/2} + \left(\prod_{j=i+1}^d t_j\right)^2\right]\\
&\leq C'\left(\prod_{j=1}^i t_j\right)\left(\prod_{j=i+1}^d t_j\right)^2
\end{align*}
\end{proof}

\end{document}
